\section{How we score a design --- and why we changed it}\label{sec:scoring}
Once a design is generated it has to be judged. We inherited Lawson's four-filter ground
truth, found that two of its filters bias selection toward rigid, well-ordered folds, and
rebuilt scoring around decomposed, epitope-specific metrics combined in a single composite.
This section follows that arc: the inherited filters, why they mislead, the decomposition that
fixes it, and the composite scorer that ranks designs and picks the best per epitope.

\subsection{Four-filter ground truth (DP3 set)}\label{sec:fourfilter}
The DP3 set itself is a \num{59}-epitope subset of the \num{1134} cleaned antibody--antigen
complexes in the AbDb directory (\texttt{configs/paths.py}, \texttt{ABDB\_CLEANED\_PDB\_DIR}).
The selection condition --- which every one of the \num{59} satisfies in \texttt{dp2.parquet}
--- is at most two epitope islands (\num{13} single-island, \num{46} dual-island) and a parent
antigen longer than \num{103} residues (the antigen lengths run \num{105}--\num{909}); the
other \num{1075} complexes are the candidates that condition drops~[verify].

A design passes when all four conditions hold: overall backbone RMSD $\le \SI{2}{\angstrom}$,
epitope-chunk RMSD $\le \SI{1}{\angstrom}$, mean predicted aligned error $< 5$, and zero
antibody residues within \SI{4}{\angstrom} of non-epitope scaffold residues
(\texttt{af3\_n\_clash\_res}). Clash detection Kabsch-aligns the AlphaFold3 epitope onto the
true epitope, then checks non-epitope scaffold heavy atoms against antibody atoms. The
antibody coordinates always come from the AbDb crystal structure, not from AlphaFold3.

We summarize a design protocol by its empirical pass rate. If a protocol produces $N$
designs $s_i$ and $f(s_i)=1$ when design $i$ clears all four filters (and $0$ otherwise),
\begin{equation}\label{eq:passrate}
\hat P_{\text{pass}} \;=\; \frac{1}{N}\sum_{i=1}^{N}\mathbf{1}\!\left[f(s_i)=1\right].
\end{equation}
This is just the fraction of designs that pass, and it is the quantity we compare between
protocols and report per epitope.

\subsection{The filters bias toward rigid folds}\label{sec:filtering}
The first thing we changed was the filtering. The original filters score a design with two
global continuous metrics, overall RMSD and mean PAE, that treat the epitope and the scaffold
as one unit. We inherited this from standard de novo practice, but it biases selection toward
globally rigid, well-ordered folds and quietly rejects designs whose scaffold is flexible while
the epitope is well positioned (or the reverse). Splitting each metric into epitope- and
scaffold-specific terms (Section~\ref{sec:decomp}) is what lets us see, and recover, those
designs.

To compare design sets on a common footing we use epitope-chunk RMSD $< \SI{2.5}{\angstrom}$
as a single bar. All three tiled 12-mer antigens pass it more often than the pooled DP3 set:

\begin{center}
\begin{tabular}{lcc}
\toprule
set & pass rate ($<\SI{2.5}{\angstrom}$) & median epitope RMSD (\si{\angstrom}) \\
\midrule
DP3 / mAb (pooled) & 30.0\% & 5.00 \\
6M0J               & 47.0\% & 2.65 \\
1D2K               & 60.2\% & 1.85 \\
4WAT               & 84.4\% & 0.43 \\
\bottomrule
\end{tabular}
\end{center}

\noindent
Even the hardest 12-mer antigen (6M0J) clears the bar more often than DP3. Figure~\ref{fig:passrates}
shows the full six-metric comparison.

\begin{figure}[h]\centering
\figorbox{dp3_vs_12mer_ungated.png}{0.95}
\caption{Ungated metric distributions (full design sets), DP3 mAb versus the three 12-mer
antigens, across epitope RMSD, global mean PAE (the four-filter metric), epitope PAE (the
decomposed metric the composite scores on), overall RMSD, pTM, AF3 clashing residues
(antibody set only), and cylinder clashes (the plain count \texttt{cylinder\_ca\_clashes}, not
the native-aware carve). Dashed line marks the median; red marks the
\SI{2.5}{\angstrom} epitope-RMSD gate with the per-set pass rate. Regenerated by
\texttt{notebooks/dp3\_vs\_12mer\_distributions.ipynb}.}
\label{fig:passrates}
\end{figure}

\paragraph{Caveats (important for interpretation).}
This comparison is not yet apples-to-apples. The 12-mer epitope is a single contiguous
12-residue stretch, while DP3 epitopes are larger and often multi-island, so 12-mer RMSD and
PAE are mechanically lower in part by definition, and the pooled DP3 30\% averages over
heterogeneous epitopes. The clean control is to stratify within DP3 by island count and size
(Section~\ref{sec:islands}). The cylinder counts are also not comparable across sets as absolute
numbers. The 12-mer medians ($\sim$31--32) run higher than DP3 ($\sim$21) mainly because the
designs are a fixed length ($\sim$104 residues) and the 12-mer epitope is small, which leaves
$\sim$91 non-epitope scaffold residues whose C$\alpha$ atoms can land in a fixed cylinder
against fewer for DP3's larger epitopes. That is a scaffold-mass effect, not more real
clashing, which is why we apply the native-aware carve and rank the cylinder per antigen
instead of against one absolute cutoff (Section~\ref{sec:noab}). As a check, the DP3 cylinder
median ($\approx 21$) tracks its real \texttt{af3\_n\_clash\_res} median ($\approx 19$).

\paragraph{Global versus epitope PAE.}
The two PAE columns show why we decompose, and they are easy to confuse. The four-filter (and
Lawson's original) uses the \emph{global} mean PAE ($<5$); on this DP3 set it sits high, a
median of $13.2$ with only $\sim$7\% of designs below 5. The scaffold region drags that value
up, since AlphaFold3 predicts it less confidently, while the epitope itself is much better
resolved (epitope PAE median $9.1$, $\sim$23\% below 5). So a design can hold its epitope well
and still fail the global PAE filter on scaffold uncertainty that need not matter for binding.
That is the bias the decomposed metrics relax, and it is why we show both columns and score
the epitope-specific PAE rather than the global one.

\subsection{Decomposed metrics}\label{sec:decomp}
The two continuous filters, overall RMSD and mean PAE, treat the epitope and the scaffold
as a single unit. We split each into epitope- and scaffold-specific terms (epitope RMSD,
scaffold RMSD, epitope PAE, scaffold PAE), retaining the zero-clash hard filter throughout.
Decomposition lets selection reward a well-placed epitope on a flexible scaffold (or the
reverse) instead of demanding global rigidity, and it lets us ask whether the optimal
rigidity profile is epitope-dependent rather than universal. Selection over the four metrics
is organized as a $2\times2$ of strict and relaxed thresholds on (epitope, scaffold), with a
flexible/flexible arm as a negative control.

\subsection{Composite scorer and per-epitope selection}\label{sec:composite}
Decomposing the metrics gives more knobs, which raised the next question: how to combine them
into one ranking and take the best designs per epitope. One config-driven scorer
(\texttt{episcaf\_analysis/score.py}) ranks designs in three steps:
(i)~transform each metric to a higher-is-better score within a scope (pooled, or per antigen);
(ii)~\textbf{weighted sum}; and (iii)~keep the \textbf{top-$k$ per group}. The scorer supports an
optional hard gate as a pre-filter, but for \emph{selection} we leave it off (see below) and let
the ranking do the work. Concretely, each design $d$ receives a score
\begin{equation}\label{eq:score}
S(d) \;=\; \sum_{m} w_m\,\varphi_m\!\bigl(x_m(d)\bigr),
\end{equation}
a weighted sum over metrics $m$ of a per-metric activation $\varphi_m$ that maps the raw
value $x_m$ to a higher-is-better score. With $\varphi$ a percentile the score is rank-based
and depends on the cohort; with a sigmoid,
\begin{equation}\label{eq:sigmoid}
\varphi(x) \;=\; \frac{1}{1+\exp\!\bigl(s\,k\,(x-x_0)\bigr)},
\qquad s=+1\ \text{if lower is better (else } -1),
\end{equation}
it is an absolute, saturating score with a tunable midpoint $x_0$ and steepness $k$. In this
form the scorer is literally a single-layer perceptron with hand-set weights (feed-forward,
no backpropagation): the activations are the per-metric nonlinearities, the $w_m$ are the
weights, and ``fitting'' it (Section~\ref{sec:open}) means learning the $w_m$ (and the
$x_0,k$) from data. The dials live in \texttt{episcaf\_analysis/presets.py}:

\begin{center}
\begin{tabular}{lc}
\toprule
metric & weight \\
\midrule
accessibility      & 0.35 \\
epitope-chunk RMSD & 0.35 \\
overall RMSD       & 0.15 \\
epitope PAE        & 0.15 \\
\bottomrule
\end{tabular}
\end{center}

\noindent
These weights are set from the DP3 binding data (Section~\ref{sec:whatpredicts}): correlating
each metric with experimental binding \emph{within} antibody, the accessibility cylinder and
epitope RMSD are the two terms that track binding (the cylinder most, and uniquely surviving the
between-antibody confound), while overall RMSD and PAE carry little signal --- so the first two
carry the weight and the latter two are kept small. This is still a hand-set prior from an
all-passing set, not a fit.

\noindent
The weights sum to one (renormalized over whichever metrics a given table actually
contains). The accessibility term depends on what the input is: with a \emph{known antibody}
(the mAb set) we use the real clash \texttt{af3\_n\_clash\_res} directly; with \emph{no antibody}
(the 12-mer set, and the sample-space search) we use the native-aware cylinder
(Section~\ref{sec:noab}), the surrogate the clash is approximated by --- the cylinder exists
precisely to stand in where there is no antibody to clash against. Either way it takes the same
accessibility weight. The 12-mer preset transforms metrics per antigen, the mAb preset pooled,
and both keep the top 5 per group. The weights are a hand-set starting point;
Section~\ref{sec:open} describes fitting them against experimental binding.

For \emph{selection} the scorer does not gate: it transforms each metric to a higher-is-better
score within its scope, weight-sums, and keeps the top designs per group, ranking the whole
population rather than hard-filtering first. A hard threshold would only discard designs the
ranking already buries, and risks dropping one that is weak on a single axis but excellent
overall; the soft per-metric penalties handle that more gracefully. (The gate remains available
as an optional pre-filter, and is used in the diagnostic statistics of Section~\ref{sec:noab},
just not in picking the designs we ship.) The scorer is deliberately simple: one metric per
input, each with
its own activation, weighted and summed, which makes it a single-layer perceptron with
hand-set weights. It is still in flux. The per-metric transform is currently a percentile
within the population, so a design's score depends on the cohort it sits in. Whether to move
to a population-independent sigmoid (an absolute, saturating per-metric score) is an open
choice we carry in Section~\ref{sec:open}.
